\chapter{SCIENTIFIC PRODUCTION}

This block focuses on the tools and practices required to document, share, and reproduce scientific work in computational science, physics, and machine learning.

\begin{center}
\textbf{Scientific production requires clear communication, version control, and reproducibility of experiments}
\end{center}

\vspace{0.5cm}

\section*{1. JUPYTER NOTEBOOKS}

\subsection*{1.1 Overview}

Jupyter Notebooks provide an interactive environment that combines code, text, and visualizations.

\subsection*{1.2 Key Features}

\begin{itemize}
\item Executable code cells
\item Markdown for documentation
\item Inline visualization
\item Step-by-step experimentation
\end{itemize}

\subsection*{1.3 Example}

\begin{minted}{python}
import numpy as np
import matplotlib.pyplot as plt

x = np.linspace(0, 10, 100)
y = np.sin(x)

plt.plot(x, y)
plt.show()
\end{minted}

\subsection*{1.4 Applications}

\begin{itemize}
\item Data analysis
\item Prototyping models
\item Educational material
\end{itemize}

\vspace{0.5cm}

\section*{2. LATEX}

\subsection*{2.1 Overview}

LaTeX is the standard typesetting system for scientific documents.

\subsection*{2.2 Mathematical Representation}

\[
E = mc^2
\]

\[
\int_a^b f(x)\,dx
\]

\subsection*{2.3 Key Features}

\begin{itemize}
\item High-quality mathematical typesetting
\item Structured documents (sections, equations, references)
\item Bibliography management
\end{itemize}

\subsection*{2.4 Applications}

\begin{itemize}
\item Academic papers
\item Reports and theses
\item Technical documentation
\end{itemize}

\vspace{0.5cm}

\section*{3. GIT AND GITHUB}

\subsection*{3.1 Version Control}

Git enables tracking changes in code and documents over time.

\subsection*{3.2 Core Concepts}

\begin{itemize}
\item Repository
\item Commit history
\item Branching and merging
\end{itemize}

\subsection*{3.3 Basic Workflow}

\begin{minted}{python}
# Conceptual commands (executed in terminal)
git init
git add .
git commit -m "Initial commit"
git push origin main
\end{minted}

\subsection*{3.4 GitHub}

\begin{itemize}
\item Remote repository hosting
\item Collaboration via pull requests
\item Issue tracking and project management
\end{itemize}

\vspace{0.5cm}

\section*{4. EXPERIMENTAL REPRODUCIBILITY}

\subsection*{4.1 Definition}

Reproducibility ensures that experiments can be repeated with consistent results.

\subsection*{4.2 Key Principles}

\begin{itemize}
\item Fixed random seeds
\item Versioned datasets
\item Controlled environments
\item Documented workflows
\end{itemize}

\subsection*{4.3 Mathematical Perspective}

Given an experiment:

\[
y = f(x, \theta)
\]

Reproducibility requires that for fixed $x$ and $\theta$, the output $y$ remains consistent.

\subsection*{4.4 Example}

\begin{minted}{python}
import numpy as np

np.random.seed(42)
data = np.random.rand(5)
print(data)
\end{minted}

\vspace{0.5cm}

\section*{FINAL CONNECTION}

Scientific production integrates:

\begin{enumerate}
\item \textbf{Exploration:} Jupyter Notebooks
\item \textbf{Documentation:} LaTeX
\item \textbf{Version Control:} Git/GitHub
\item \textbf{Reliability:} Reproducibility practices
\end{enumerate}

\vspace{0.5cm}

\section*{STUDY ROADMAP}

\subsection*{Phase 1 --- Notebooks}
\begin{itemize}
\item Jupyter basics
\item Combining code and text
\end{itemize}

\subsection*{Phase 2 --- Documentation}
\begin{itemize}
\item LaTeX structure
\item Mathematical writing
\end{itemize}

\subsection*{Phase 3 --- Version Control}
\begin{itemize}
\item Git fundamentals
\item GitHub collaboration
\end{itemize}

\subsection*{Phase 4 --- Reproducibility}
\begin{itemize}
\item Experiment tracking
\item Environment management
\end{itemize}

\subsection*{Phase 5 --- Integration}
\begin{itemize}
\item End-to-end scientific workflows
\item Research project organization
\end{itemize}